% Physical unit 013; accepted Session 02.
\GordanUnitStart{013}{printed page 8}
\GordanTopFolio{8}
\begin{center}
{\large CAPUT III.}\\[0.8em]
\S\ 1.
\end{center}

Nominetur:
\[
 a\int_0^{\frac\pi2} d\varphi\sqrt{1-e^2+e^2\sin^2 l_0\cos^2\varphi}
\]
periodus extensionis et
\[
 \int_0^{\frac\pi2}\cos l_0\,d\varphi\,
 {\sqrt{1-e^2+e^2\sin^2 l_0\cos^2\varphi}\over 1-\sin^2 l_0\cos^2\varphi}
\]
periodus longitudinis lineae geodeticae $l_0$, quae, cum pendeant tantum ab angulo $l_0$, significentur $\Sigma(l_0)$ et $\mathcal A(l_0)$. Quoniam harum periodorum valores deinceps gravissimi momenti erunt, jam series evolvendae sunt, quibus illi facile quam accuratissime computentur.

\begin{center}\S\ 2.\end{center}

Ad inveniendum $\Sigma(l_0)$ sub signo integrationis
\[
 \sqrt{1-e^2+e^2\sin^2 l_0\cos^2\varphi}
\]
secundum theorema binomicum evolvimus, ubi, ut breves simus, ponamus:
\[
 {1\over4}\,{e^2\over1-e^2}={\varepsilon\over(1-\varepsilon)^2}=o
\]
sequitur jam:
\[
\begin{aligned}
 {1\over b}\Sigma(l_0)
 &=\int_0^{\frac\pi2}d\varphi\{1+2o\sin^2l_0\cos^2\varphi-2o^2\sin^4l_0\cos^4\varphi\\
 &\qquad{}+4o^3\sin^6l_0\cos^6\varphi-10o^4\sin^8l_0\cos^8\varphi\ldots\}\\
 &=\int_0^{\frac\pi2}d\varphi+2o\sin^2l_0\int_0^{\frac\pi2}\cos^2\varphi\,d\varphi\\
 &\quad{}-2o^2\sin^4l_0\int_0^{\frac\pi2}\cos^4\varphi\,d\varphi
 +4o^3\sin^6l_0\int_0^{\frac\pi2}\cos^6\varphi\,d\varphi\\
 &\quad{}-10o^4\sin^8l_0\int_0^{\frac\pi2}\cos^8\varphi\,d\varphi\ldots
\end{aligned}
\]
Singula integralia formam praestant hanc:
\[
 \int_0^{\frac\pi2}\cos^{2m}\varphi\,d\varphi
\]
